% Project:  TeX PhD Motivation Letter Template
% Author:   Deep Ghuge

% EDIT main.tex directly

\documentclass[11pt]{report}
\usepackage[margin=0.85in]{geometry}
\usepackage{graphicx}

\usepackage{hyperref}

\newcommand{\mediumsize}{\fontsize{15pt}{16pt}\selectfont}

\begin{document}
	\begin{titlepage}
		
		\begin{minipage}[t]{0.90\textwidth}
			\hfill
			\raggedleft
			\textbf{Edwin Varghese} \\
			London, United Kingdom \\
			\today\\
			\href{mailto:edwinshajivarghese@gmail.com}{edwinshajivarghese@gmail.com} \\
			\url{www.linkedin.com/in/edwinvarghese1}
		\end{minipage}
		
		\raggedright \textbf{Dr Christos Anastopoulos} \\ School of Mathematical and Physical Sciences \\ University of Sheffield \\ Sheffield, UK\\ \href{mailto:c.anastopoulos@sheffield.ac.uk}{c.anastopoulos@sheffield.ac.uk}\\
		
		\vspace{0.7em}
		
		\raggedright Dear Dr Christos Anastopoulos,\\
		\vspace{0.7em}
		
		I am a fourth-year MSci Physics student at Royal Holloway, University of London, and I am writing to apply for the PhD position to conduct electroweak measurements with the ATLAS experiment. I am eager to work with the new 13.6 TeV Run 3 dataset and contribute to precision SM measurements, which I see as the next logical step for my development as an experimental particle physicist.\\
		\vspace{0.7em}
		
		My experience as a third-year student working on a project exploring the phenomenological structure of Higgs-boson decays convinced me that experimental collider physics is the field I want to work in. Driven by curiosity, I learnt that the introduction of the Higgs field and the Higgs mechanism explains how particles acquire mass without breaking the gauge structure of the Standard Model and also makes precise, testable predictions that experiments like ATLAS can measure. Yet, the Higgs does not explain why the Yukawa couplings have the values they do, nor does it account for the hierarchy of fermion masses, or the absence of additional scalar states. With this theoretical motivation in place, I could not wait to explore how such ideas are pursued in practice, and it was through research placements that my interests were justified.\\
		\vspace{0.7em}
		
		During a six-week research placement at my university, I implemented and trained neural-network-based classifiers to evaluate a novel MCMC convergence metric ($R^*$) and applying the method to simulated datasets from the T2K and NO$\nu$A experiments. I benchmarked $R^*$ against the industry standard Gelman-Rubin ($\hat{R}$) statistic by measuring the metric’s sensitivity to statistical noise, showing that $R^*$ detects smaller deviations between chains and provides complementary global convergence information where  $\hat{R}$ is insensitive. Over the course of the project I developed practical skills in applying and validating ML methods for physics analyses and I am eager to apply and further develop these skills as part of an analysis team making precise electroweak measurements.\\
		\vspace{0.7em}
		
		Complementing this, my UKAEA plasma-simulation internship focused on integrating and benchmarking simulation codes for MAST-Upgrade plasmas. I automated the generation of RABBIT input files --- a novel, fast, neutral-beam injection code for plasma simulation --- containing data about the initial state of the plasma, and the plotting of post-simulation outputs with Python. RABBIT was benchmarked against the higher-fidelity NUBEAM/TRANSP modelling chain using the OMFIT workflow manager. Through this internship I learnt how to organise and check large sets of simulation outputs in a systematic way. I also gained experience writing Python scripts with good documentation to ensure the work could be reproduced and shared easily. Having built confidence working with complex simulation tasks, I then undertook projects directly relevant to collider phenomenology at Royal Holloway.\\
		\vspace{0.7em}
		
		My third-year project explored how variations in the Higgs boson mass affect decay channel contributions; I solved formulae numerically using Python and produced logarithmic plots of branching ratios and widths, presenting my results to faculty and peers. I also completed a research review of Higgs searches at the LHC, in which I discussed the relevant production and decay channels, and the impact of higher-order QCD and electroweak corrections on discovery sensitivity. My ongoing Master’s project involves deriving the hadronic W production cross-section, combining analytic phase-space reduction with numerical integration and Monte-Carlo sampling to produce differential distributions. The analysis, ML and simulation skills would translate naturally to precision Higgs and electroweak measurements.\\
		\vspace{0.7em}
		
		Beyond technical skills, I have leadership and communication experience as Academic Officer of the RHUL Physics Society, where I ran LaTeX and Python workshops, organised a weekly PhD seminar series, and coordinated student-led conferences. I enjoy teaching, presenting, and collaborating --- qualities that I believe fit well with the collaborative ethos of ATLAS.\\
		\vspace{0.7em} 
		
		My experience across machine learning, simulation work and collider-focused projects has convinced me that I want to build my career in experimental particle physics. My responsibilities in collaborative research environments have taught me the importance of clear documentation, careful validation and open communication. I am eager to develop and consolidate my skills as an analyst in a new setting where I can work with real ATLAS data, and it will be a pleasure to exchange ideas with and learn from researchers who are leading precision Higgs and electroweak measurements within ATLAS. What draws me most to this project is the opportunity to work on analyses that build directly on the electroweak and Higgs theory I have studied, while gaining experience with ATLAS datasets and the statistical tools used in current Run 3 measurements. I am excited about the prospect of contributing to a search where new physics may genuinely emerge and I look forward to growing into an independent researcher within the ATLAS community.\\
		\vspace{0.7em}
		
		
		Thank you for considering my application.\\
		\vspace{5em}
		
		\raggedright Yours sincerely,\\
		\vspace{0.7em}
		\textbf{Edwin Varghese}
		
	\end{titlepage}
\end{document}